Table~\ref{tab:cross} and Fig.~\ref{fig:cross} report the \numCross{} unseen plants of the box, none of which was seen in training and none of which is known to any controller, with a front and a back transition each. HOT reaches the target band within the \SI{20}{s} window on \hotCrossFrontReached{} of the \hotCrossFrontN{} front and \hotCrossBackReached{} of the \hotCrossBackN{} back transitions, in a median of \hotCrossFrontTmed{}~s and \hotCrossBackTmed{}~s when it does, keeps the lift fraction above \hotCrossLf{} on every plant and violates no constraint on any of them; the module intervened at \hotCrossFrontIntervPct{}\% of the front and \hotCrossBackIntervPct{}\% of the back updates. The Wilcoxon test on the time to target is significant at $p<0.05$ in favour of HOT against every controller except its own unfiltered version, and the Friedman mean rank over the \numFriedmanN{} transitions flown by all controllers is \hotCrossRank{} for HOT against \bestOtherCrossRank{} for the best alternative, \bestOtherCrossName{}. The price of the guarantee is energy, about one watt-hour more than the classical controllers, spent by the rotors during the shorter transition, and the lift fraction, which the schedule and the gain-scheduled PID hold at 0.97 by taking twice the time. The transitions that miss the window fall into two groups. In the front transition the policy commands a nose-down attitude on strong-wing plants that cannot be certified against the light-wing plants still retained in the set, so the certified command holds the rotors and the aircraft climbs instead of accelerating; in the back transition \hotCrossBackStalledUnver{} plants stall after a single unverified backup step left the policy in a state it had not learned from, so it decelerated by drag alone and stopped just above the band, and \hotCrossBackStalledNoInt{} stall without any intervention. The transformer without the module, HOT$-$LSM, reaches \hotnfCrossReached{} of the transitions in \hotnfCrossT{}~s but violates a constraint on \hotnfCrossViol{}\% of the samples over the family, the threshold rule on \thrCrossViol{}\%, the schedule on \schCrossViol{}\% and PPO on \ppoCrossViol{}\%; the gain-scheduled PID and the MPC violate nothing but take \pidCrossT{} and \mpcCrossT{}~s, and the MPC reaches \mpcCrossReached{} of its transitions. The module is therefore what turns the fastest policy into a safe one, at the cost of the stalls just described.
\begin{table*}[t]\centering\footnotesize
\caption{Generalisation across the plant family: \numCross{} unseen plants, front and back transitions, mean $\pm$ standard deviation; W: Wilcoxon signed-rank test against HOT ($+$: HOT better, $=$: no difference, $-$: HOT worse, $p<0.05$); rank: Friedman mean rank on the time to target (lower is better).}\label{tab:cross}
\begin{tabular}{l|c|cc|cc|cc|c|c}\toprule
controller & reached & time [s] & W & energy [Wh] & W & min $\lf$ & W & viol.\ [\%] & rank\\\midrule
THR & 0.96 & 8.0$\pm$3.9 & $+$ & 3.18$\pm$0.67 & $-$ & 0.88$\pm$0.06 & $+$ & 1.6 & 3.82\\
SCH & 0.87 & 9.4$\pm$2.9 & $+$ & 3.62$\pm$0.44 & $-$ & 0.97$\pm$0.03 & $-$ & 2.5 & 5.24\\
GS-PID & 0.95 & 9.7$\pm$4.5 & $+$ & 3.50$\pm$0.34 & $-$ & 0.97$\pm$0.03 & $-$ & 0.0 & 5.06\\
MPC & 0.50 & 15.1$\pm$3.4 & $+$ & 3.44$\pm$0.22 & $-$ & 0.96$\pm$0.04 & $-$ & 0.0 & 7.28\\
HOT$-$LSM & 0.98 & 5.2$\pm$2.1 & $-$ & 4.39$\pm$0.28 & $=$ & 0.92$\pm$0.07 & $=$ & 1.6 & 1.89\\
PPO & 0.50 & 8.3$\pm$1.1 & $+$ & 6.48$\pm$2.49 & $+$ & 0.60$\pm$0.31 & $+$ & 36.6 & 5.61\\
SAC & 0.87 & 5.9$\pm$2.0 & $+$ & 3.66$\pm$0.67 & $-$ & 0.91$\pm$0.03 & $+$ & 0.3 & 3.86\\
\textbf{HOT} & 0.89 & 5.5$\pm$2.5 &  & 4.38$\pm$0.25 &  & 0.94$\pm$0.02 &  & 0.0 & 3.24\\
\bottomrule\end{tabular}
\end{table*}
\begin{figure*}[t]\centering\includegraphics[width=\textwidth]{fig07_cross.pdf}
\caption{Distribution over the \numCross{} unseen plants of the time to target, the energy, the minimum lift fraction and the maximum angle of attack for each controller (boxes: quartiles; whiskers: extremes; dotted: bounds).}\label{fig:cross}
\end{figure*}
